% !TeX program = lualatex
% !TeX encoding = utf8
% !TeX spellcheck = uk_UA
% !TeX root =../ClassicalEletrodynamics.tex

%=========================================================
\chapter{Варіаційний принцип для рівнянь електродинаміки}\label{\currfilebase}
%=========================================================

Рівняння динаміки частинок та полів можна пов'язати з певними варіаційними принципами, які широко використовують в теоретичній фізиці. Ці принципи пов'язані з поняттями дії та лагранжіану. В цьому розділі варіаційний принцип найменшої дії проілюстровано на прикладі рівнянь електродинаміки. Буде подано функціонали дії, з яких можна отримати рівняння руху заряджених частинок і рівняння електромагнітного поля. Цей розділ спирається на знання основ варіаційного числення.

%% --------------------------------------------------------
\section{Заряд у зовнішньому електромагнітному полі}
%% --------------------------------------------------------

Розглянемо траєкторію (світову лінію) точкової частинки $x^\mu = x^\mu(p)$, де $p$ --- деякий неособливий параметр. На ділянці траєкторії між точками $a = x^\mu(p_1)$ та $b= x^\mu(p_2)$ розглянемо інтеграл
\begin{equation}\label{eq4.1}
S_0 = -mc \int_a^b ds = -mc \int_{p_1}^{p_2} \sqrt{\eta_{\mu\nu} \frac{dx^\mu}{dp} \frac{dx^\nu}{dp}} \, dp
\end{equation}
що має назву дії для вільної частинки; вибір коефіцієнтів забезпечує відповідність з нерелятивістським наближенням для дії. Також розглянемо інтеграл
\begin{equation}\label{eq4.2}
S_{\text{int}} = -\frac{q}{c} \int_a^b A_\mu dx^\mu = -\frac{q}{c} \int_a^b A_\mu \frac{dx^\mu}{dp} \, dp
\end{equation}
що описує взаємодію частинки з електромагнітним полем; тут $A_\mu$ --- чотиривектор потенціалу, $q$ --- заряд частинки. Легко перевірити, що параметр $p$ можна вибрати довільно; при будь-якій заміні $p \to p' = p'(p)$ вигляд~\eqref{eq4.1} та~\eqref{eq4.2} не змінюється. Уведемо дію для частинки у зовнішньому полі
\begin{equation*}
S = S_0 +S_{\text{int}}.
\end{equation*}

Дія $S$ є функціоналом, що зіставляє кожній траєкторії певне число. Далі ми розглянемо першу варіацію $\delta S$ по траєкторії частинки. Зауважимо, що, як буде видно далі, повний функціонал дії є сумою дій для окремих частинок плюс дія для поля. Але при обчисленні варіацій $\delta S$ траєкторії різних частинок вважаються незалежними, тому на цьому етапі досить розглянути дію для однієї частинки. Дія для поля не залежить від змінних частинок і не впливає на варіацію $\delta S$ при варіаціях траєкторії частинки $x_0^\mu(p) \rightarrow x_0^\mu(p) + \delta x^\mu(p).
$.

Варіаційний принцип найменшої дії $\delta S = 0$ для частинки у зовнішньому полі можна сформулювати так:
\begin{postulate}{}
Перша варіація $\delta S = \delta S_0 +\delta S_{\text{int}}$ дорівнює нулю на траєкторії $x_0^\mu(p)$, тоді й тільки тоді, якщо ця траєкторія задовольняє рівнянням руху. Траєкторію $x_0^\mu(p)$ називають опорною.
\end{postulate}

Зауважимо, що, як відомо з варіаційного числення, умова $\delta S = 0$ є лише необхідною умовою мінімуму функціоналу дії. Аналіз достатніх умов мінімуму потребує додаткового розгляду.

Покажемо, що наслідком принципу найменшої дії згідно до формул~\eqref{eq4.1},~\eqref{eq4.2} є релятивістські рівняння руху частинки в електромагнітному полі. Розглянемо траєкторії в околі опорної світової лінії $x^\mu(p)$. При малих змінах траєкторії $\delta x^\mu(p)$ числове значення функціоналу $S_p$ змінюється. Нас цікавитимуть малі зміни дії
\begin{equation*}
\delta S_p = S_p[x^\mu + \delta x^\mu] - S_p[x^\mu],
\end{equation*}
при довільних $\delta x^\mu(p)$ за стандартної умови
\begin{equation}\label{eq4.3}
\delta x^\mu(p_a) = \delta x^\mu(p_b) = 0
\end{equation}
(так звана варіація з закріпленими кінцями). Малі зміни дії розглядаємо з точністю до лінійних по $\delta x^\mu$ членів, тобто розгляд обмежується першою варіацією. Дію $S_p$ запишемо у вигляді
\begin{equation*}
S_p = \int_a^b L(x^\mu, \dot{x}^\mu) dp,
\end{equation*}
де $L = -mc\sqrt{g_{\mu\nu}\dot{x}^\mu \dot{x}^\nu} - \frac{q}{c} A_\mu \dot{x}^\mu$; тут позначено $\dot{x}^\mu = \frac{dx^\mu}{dp}$.

Покажемо, що з умови $\delta S = 0$ випливають рівняння руху зарядженої частинки
\begin{equation}\label{eq4.4}
mc \frac{du_\mu}{ds} = \frac{q}{c} F_{\mu\nu} u^\nu, \text{ де } u^\mu = \frac{dx^\mu}{ds}.
\end{equation}
За стандартною процедурою варіаційного числення умова $\delta S = 0$ з урахуванням~\eqref{eq4.3} приводить до рівнянь Ейлера-Лагранжа для функції $L$:
\begin{equation*}
\frac{\partial L}{\partial x^\mu} - \frac{d}{dp} \left( \frac{\partial L}{\partial \dot{x}^\mu} \right) = 0,
\end{equation*}
або, обчислюючи похідні від $L$,
\begin{equation*}
-\frac{q}{c} \frac{\partial A_\nu}{\partial x^\mu} \dot{x}^\nu + \frac{d}{dp} \left( mc \frac{\dot{x}_\mu}{\sqrt{\dot{x}_a \dot{x}^a}} + \frac{q}{c} A_\mu \right) = 0,
\end{equation*}
де $\dot{x}_\mu = g_{\mu\nu} \dot{x}^\nu$. Перейдемо від параметра $p$ до $s$:
\begin{equation*}
mc \frac{d}{ds} \frac{dx_\mu}{ds} + \frac{q}{c} \frac{d A_\mu}{ds} - \frac{q}{c} \frac{\partial A_\nu}{\partial x^\mu} \frac{dx^\nu}{ds} = 0.
\end{equation*}
Розкриваємо похідну в лівій частині:
\begin{equation*}
mc \frac{du_\mu}{ds} + \frac{q}{c} \left( \frac{\partial A_\mu}{\partial x^\nu} \frac{dx^\nu}{ds} - \frac{\partial A_\nu}{\partial x^\mu} \frac{dx^\nu}{ds} \right) = 0.
\end{equation*}
Результатом є рівняння, які збігаються з~\eqref{eq4.4} після підняття індексу $\mu$. Звернемо увагу, що хоча в лагранжіані $L$ фігурує 4-вектор потенціал $A_\mu$, він з'являється в явному виді в рівняннях руху лише через тензор електромагнітного поля $F_{\mu\nu}$, який, таким чином, і є спостережуваною величиною.

\begin{problem}
Релятивістська заряджена частинка з масою $m$ та зарядом $q$ рухається у полі нерухомого заряду $Q$. Отримати рівняння для тривимірної траєкторії частинки.
\end{problem}

\textit{Розв'язання.} Коваріантний 4-вектор-потенціал нерухомого заряду в початку координат є $A_\mu = (\phi, 0, 0, 0) = (Q/r, 0, 0, 0)$. Якщо у виразі для дії перейти до параметру $t$, підінтегральний вираз дає тривимірну функцію Лагранжа зарядженої частинки у зовнішньому полі
\begin{equation*}
L = -mc^2 \sqrt{1 - v^2/c^2} - q\phi, \quad \phi = Q/r.
\end{equation*}
Звідси маємо інтеграл енергії
\begin{equation}\label{eqR1}
\mathcal{E} = \frac{mc^2}{\sqrt{1 - v^2/c^2}} + \frac{qQ}{r} = \text{const}.
\end{equation}
В сферичних координатах в площині $\theta = \pi/2$ інтеграл моменту є
\begin{equation}\label{eqR2}
M = \frac{mr^2 \dot{\varphi}}{\sqrt{1 - v^2/c^2}} = \text{const}.
\end{equation}
Комбінуючи~\eqref{eqR1} та~\eqref{eqR2}, позбавимось $dt$ та отримаємо
\begin{equation}\label{eqR3}
\left( \frac{dr}{d\varphi} \right)^2 + r^2 = \frac{r^4}{M^2 c^2} \left[ \left( \mathcal{E} - \frac{qQ}{r} \right)^2 - m^2 c^4 \right].
\end{equation}
З урахуванням цього співвідношення, позначаючи $u = 1/r$, $\alpha = qQ/Mc$,
\begin{equation*}
\left( \frac{du}{d\varphi} \right)^2 + u^2 = \frac{\mathcal{E}^2 - m^2 c^4}{M^2 c^2} - \frac{2\mathcal{E} qQ}{M^2 c^2} u + \frac{q^2 Q^2}{M^2 c^2} u^2.
\end{equation*}
Звідси та з~\eqref{eqR1} маємо
\begin{equation*}
\frac{d^2 u}{d\varphi^2} + u (1 - \alpha^2) = -\frac{\mathcal{E} qQ}{M^2 c^2}
\end{equation*}
або
\begin{equation*}
\frac{d^2}{d\varphi^2} \left( u + \frac{\mathcal{E} qQ}{M^2 c^2 (1 - \alpha^2)} \right) + u (1 - \alpha^2) = 0.
\end{equation*}
За $1 - \alpha^2 > 0$ рівняння має вигляд закону збереження енергії для осцилятора зі зміщеним центром; його розв'язками є тригонометричні функції. За $1 - \alpha^2 < 0$ розв'язками є гіперболічні функції, зокрема, можливе падіння на центр $r \to 0$.

%% --------------------------------------------------------
\section{Дія для електромагнітного поля}
%% --------------------------------------------------------

Мета цього підрозділу --- сформулювати варіаційний принцип, з якого випливають рівняння електромагнітного поля Максвелла. Згідно з сучасними уявленнями, дія для системи “поле + джерела” має складатися з суми функціоналів дії для поля, для зарядів та доданка, що описує взаємодію поля із зарядами:
\begin{equation}\label{eq4.5}
S = S_f + S_m + S_{mf}.
\end{equation}
Оскільки поля задано в усьому просторі, функціонал дії повинен містити інтегрування по чотиривимірному об'єму, а не тільки по траєкторіях, як у випадку частинок. Тому подаватимемо $S$ у вигляді такого інтегралу:
\begin{equation}\label{eq4.6}
S = \int_{\Omega} \Lambda dx^4.
\end{equation}
Незалежними змінними, відносно яких обчислюватиметься варіація, вважаємо компоненти 4-вектора потенціалу $A_\mu$. Область інтегрування $\Omega$ може охоплювати весь простір, а може бути скінченною. Варіацію функціоналу дії розглядаємо за умови, що зміни польових функцій обертаються на нуль на межі області інтегрування $\partial\Omega$.

\textit{Зауваження:} Якщо змінити підінтегральний вираз в~\eqref{eq4.6}, додавши до нього повну 4-дивергенцію $\partial_\mu F^\mu$, де $F^\mu$ --- будь-яка функція від польових змінних (але не від їх похідних), то при обчисленні варіації додається вираз
\begin{equation*}
\delta \int \partial_\mu F^\mu dx^4 = \int \partial_\mu (\delta F^\mu) dx^4,
\end{equation*}
що обертається на нуль в силу граничних умов на $\partial\Omega$ після застосування формули Остроградського-Гаусса. Тому додавання повної 4-дивергенції до $\Lambda$ не впливає на рівняння, що випливають з варіаційного принципу.

Вирази для $S_m$ та $S_{mf}$ отримаємо з~\eqref{eq4.1},~\eqref{eq4.2} за допомогою підсумовування дій для окремих частинок.
\begin{equation}\label{eq4.7}
S_m = -\sum mc \int ds, \quad S_{mf} = -\sum \frac{q}{c} \int A_\mu dx^\mu.
\end{equation}
Як було зазначено, ці вирази слід подати у вигляді інтегралів. Зробимо це лише для $S_{mf}$, оскільки $S_m$ не впливає на виведення рівнянь поля при обчисленні варіації $\delta S$ по $A_\mu$. Для неперервного розподілу ми зіставляємо кожному елементу заряду $dq$ його траєкторію $x^\mu(s)$; відповідно
\begin{equation*}
S_{mf} = -\int \frac{1}{c} A_\mu \frac{dx^\mu}{dt} dq dt,
\end{equation*}
де $dq = \rho dV$ --- об'ємна густина заряду, $\rho$, $j^\mu = \rho \frac{dx^\mu}{dt}$ --- 4-вектор густини струму. Тоді, переходячи в $S_{mf}$ до інтегралу по області $\Omega$ чотиривимірного простору маємо:
\begin{equation}\label{eq4.8}
S_{mf} = -\frac{1}{c^2} \int_{\Omega} A_\mu j^\mu dx^4.
\end{equation}

Виникає питання, чи не втратимо ми калібрувальну інваріантність теорії завдяки тому, що в~\eqref{eq4.8} явно входить 4-потенціал? Відсутність тут якихось непорозумінь випливає в першу чергу з того, що з принципу найменшої дії буде отримано правильні рівняння поля, що є калібрувально інваріантними. Але навіть на даному етапі можна пересвідчитися, що вираз~\eqref{eq4.8} є калібрувально інваріантним, якщо врахувати закон збереження заряду. Покажемо, що~\eqref{eq4.8} не змінюється при калібрувальних перетвореннях
\begin{equation*}
A_\mu \to A_\mu + \partial_\mu \psi,
\end{equation*}
де $\psi = 0$ на межі області $\partial\Omega$. При такому перетворенні інтеграл~\eqref{eq4.8} отримує доданок
\begin{equation*}
\delta_\psi S_{mf} = -\frac{1}{c^2} \int j^\mu \partial_\mu \psi dx^4 = -\frac{1}{c^2} \int \partial_\mu (j^\mu \psi) dx^4 + \frac{1}{c^2} \int \psi \partial_\mu j^\mu dx^4.
\end{equation*}
Інтеграл від дивергенції зникає на межі $\partial\Omega$, а другий доданок дорівнює нулю завдяки закону збереження заряду $\partial_\mu j^\mu = 0$. Це ілюструє зв'язок калібрувальної інваріантності із збереженням заряду.

Підберемо вираз для $\Lambda_f$, що приводить до рівнянь класичної електродинаміки. Поле для пошуків $\Lambda_f$ значно звужується, якщо врахувати такі обставини. По-перше, рівняння електромагнітного поля мають однаковий вигляд в усіх інерціальних системах відліку. Тому доцільно конструювати $\Lambda_f$ як скаляр за допомогою згорток вектора $A_\mu$ та його похідних $\partial_\nu A_\mu$. По-друге, в рівняннях руху фігурує не безпосередньо $A_\mu$, а величини $F_{\mu\nu} = \partial_\mu A_\nu - \partial_\nu A_\mu$. Саме $F_{\mu\nu}$ мають фізичний зміст як величини, які можливо вимірювати, причому вони інваріантні відносно калібрувальних перетворень. Щоб зберегти калібрувальну інваріантність в рівняннях поля, будемо конструювати $\Lambda_f$ саме з величин $F_{\mu\nu}$. Ці величини задовольняють рівнянням~\eqref{eq3.2}, які запишемо у вигляді
\begin{equation}\label{eq4.9}
\partial_\lambda F_{\mu\nu} + \partial_\mu F_{\nu\lambda} + \partial_\nu F_{\lambda\mu} = 0.
\end{equation}
Щоб отримати лінійні рівняння поля, природно обмежитися лише квадратичними комбінаціями $F_{\mu\nu}$, які є скалярами відносно групи Лоренца. Таких комбінацій, з точністю до сталого множника, лише дві:
\begin{equation*}
I_1 = F_{\mu\nu} F^{\mu\nu} \quad \text{ та } \quad I_2 = F_{\mu\nu} \tilde{F}^{\mu\nu}.
\end{equation*}
Точніше, $I_2$ --- псевдоскаляр; він змінює знак при інверсіях координат; він зводиться до повної дивергенції:
\begin{equation*}
F_{\mu\nu} \tilde{F}^{\mu\nu} = \partial_\mu (2\epsilon^{\mu\nu\alpha\beta} A_\nu \partial_\alpha A_\beta),
\end{equation*}
де було враховано антисиметрію $\epsilon^{\mu\nu\alpha\beta}$ та рівняння~\eqref{eq4.9}, що мають місце завдяки вигляду $F_{\mu\nu}$ за побудовою. Інтеграл від $I_2$ по області $\Omega$ зводиться до інтегралу по межі $\partial\Omega$ і не дає внеску в рівняння поля всередині $\Omega$. Таким чином, в нашому розпорядженні залишається лише інваріант $I_1$, тому можна вибрати
\begin{equation}\label{eq4.10}
\Lambda_f = -\frac{1}{16\pi} F_{\mu\nu} F^{\mu\nu}.
\end{equation}
Вибір сталого множника зроблено заздалегідь так, щоб забезпечити відповідність рівнянням Максвелла в гаусовій системі одиниць. Можна показати, що вибір іншого множника вплине лише на визначення одиниці заряду. Враховуючи внесок взаємодії поля і зарядів, покладемо
\begin{equation}\label{eq4.11}
S = \int_{\Omega} \left( -\frac{1}{16\pi} F_{\mu\nu} F^{\mu\nu} - \frac{1}{c^2} A_\mu j^\mu \right) dx^4,
\end{equation}
де $F_{\mu\nu} = \partial_\mu A_\nu - \partial_\nu A_\mu$. Нагадаємо, що незалежними польовими змінними в~\eqref{eq4.11} вважаємо $A_\mu$. Обчислимо $\delta S$ за умови, що варіації польових змінних на межі області інтегрування зникають:

Маємо
\begin{equation*}
\delta(F_{\mu\nu} F^{\mu\nu}) = 2 F^{\mu\nu} \delta F_{\mu\nu} = 2 F^{\mu\nu} (\partial_\mu \delta A_\nu - \partial_\nu \delta A_\mu) = 4 F^{\mu\nu} \partial_\mu \delta A_\nu;
\end{equation*}
\begin{equation*}
\int F^{\mu\nu} \partial_\mu \delta A_\nu dx^4 = \int \partial_\mu (F^{\mu\nu} \delta A_\nu) dx^4 - \int \delta A_\nu \partial_\mu F^{\mu\nu} dx^4 = - \int \delta A_\nu \partial_\mu F^{\mu\nu} dx^4,
\end{equation*}
де застосовано формулу Остроградського-Гаусса і враховано умову на $\partial\Omega$. Звідси
\begin{equation*}
\delta S = \int \left( \frac{1}{4\pi} \partial_\mu F^{\mu\nu} - \frac{1}{c^2} j^\nu \right) \delta A_\nu dx^4.
\end{equation*}
Згідно з принципом стаціонарної дії, цей вираз має дорівнювати нулю для довільних варіацій $\delta A_\nu$ всередині $\Omega$. Це можливо лише коли вираз в фігурних дужках дорівнює нулю: $\frac{1}{4\pi} \partial_\mu F^{\mu\nu} - \frac{1}{c^2} j^\nu = 0$, або остаточно
\begin{equation}\label{eq4.12}
\partial_\mu F^{\mu\nu} = \frac{4\pi}{c^2} j^\nu.
\end{equation}
Разом із~\eqref{eq4.9} це рівняння утворює повну систему рівнянь Максвелла~\eqref{eq3.1},~\eqref{eq3.2} для тензора електромагнітного поля $F_{\mu\nu}$.